Moeder noemde als leus
haar vaak "eigenwijsneus!".

Emmy stak haar neus in
alles wat haar aanging

en wist het dan beter.
Dat mocht wel discreter.

“Weet jij als geleerd wicht,
waar venusheuvel ligt?”

plaagde moeder miskend.
“Ik ben niet zo bekend

met de landstreek, Mama!”
Haar zus grapte dan na.

Die mentale groeipijn
vergrootte haar welzijn.

Het klooster van de stad
lag op haar vrome pad.

Buiten de richtlijnen
werd daar door begijnen,

als leken-lerares,
met wisselend succes

onderwijs gegeven.
Ze leerde daar lezen

en liet zich bekeren.
Emmy hield van leren,

maar een boerenleven
was voorts niet haar streven.

Dat leven vond ze saai
en bovendien te taai.

Het klooster stond net los
van heideveld en bos

op een verhoogd terras
omringd door een moeras.

Deze afzondering
gaf elke kloosterling

de aangewezen rust
en koos men dus bewust.

Het kerkepad was smal,
onverhard met houtwal.

De groene lindelaan
liet de bezoeker gaan

richting een ophaalbrug.
Die hing meestal terug.

Houten palissaden
voor dieren en kwaden

lagen binnen de gracht.
Een portier hield de wacht.

Het gebouw van baksteen
was in het algemeen

voor activiteiten
en festiviteiten

van religieuze aard
en gold als huis en haard.

De cisterciënzer
overdaadbegrenzer

hield de gotische stijl
binnen het levenspeil

van de gewone man,
maar met genoeg elan.

Bim-bam kon men horen
uit de klokkentoren

voor een ferme oproep
om met de kloostergroep

in elkanders midden
te werken of bidden.

De klus werd trouw geklaard
in de moestuin, boomgaard,

kapel, scriptorium
of refectorium.

De meest wijze nonnen
waren voor haar bronnen

vol met inspiratie.
Op een dag vroeg Emmy

straf aan een kloostermeid;
"Wat is barmhartigheid?”

Haar verrassend antwoord
was nogal ongehoord;

“Mocht jij dat ontdekken
wil jij het verstrekken

aan mij als kloosterling.
Want meelij is het ding

wat ik hier constant zoek
sinds mijn eerste bezoek.

Mijn werkmotivatie,
mijn geachte Emmy,

is een verkooppraatje
om in een goed blaadje

te komen bij de heer.
Niets minder en niets meer.”

Intuïtie als spil.
Geloof uit vrije wil.

Niet gegrond in scholing,
maar in bewustwording.

Als George Fox' innerlicht,
dat zich puur op God richt.

De abdis hoorde het
en zag dit als een smet

op het kloosterleven.
De non stond te beven
tijdens de uitbrander.
Ze was geen voorstander

van de ontaarde leer,
die ze stram keer op keer

luidkeels moest aanhoren.
Emmy kon zich storen

aan de oude Abdis.
Haar scherpste ergernis

kwam uit haar onderbuik.
De zweem van machtsmisbruik

middels schijnheiligheid.
Familiariteit

uit haar allerdiepste.
Al trotserend riep ze,

hoezeer ook de kriegel
van een mens als spiegel,

tevens vertrouwdheid op.
“Hoe kon ooit de bisschop

dit intens valse kreng,
meedogenloos en streng,

hier hebben aangesteld?”
Een non had haar verteld

dat de abdis voorheen
als apostel verscheen

van een nieuwe profeet
voor wie ze missies deed.

Ze noemde zich ‘Messing’
als vreemde aanduiding.

Ze was toen zelfs getrouwd.
Ze werd als smiecht beschouwd

door de kloosterlingen.
Ze moest zich opdringen

en nam vele horden
om gehoord te worden.

Dat deed slechts de bisschop.
Zo kwam ze aan de top.

Als ze op haar kop kreeg,
omdat haar leer uitsteeg

boven die van de kerk,
dan ging ze aan het werk

met nog meer toewijding.
De kloeke kloosterling,

die haar had verraden,
kon zijn toekomst raden.

Haar gemoed werd driftig.
Haar adem werd giftig.

Wie haar kop eraf sloeg,
en zo om revanche vroeg,

die zag er twee daarna.
Net als bij de Hydra.

Behalve de abdis
gaf men Emmy kennis

die ze buiten niet kreeg.
Haar nieuwsgierigheid steeg.

Een spiritueel pad
leek haar daarom wel wat.

Ze duidde deze weg,
passend bij haar aanleg

voor het reflecteren,
leren, mediteren

en het contempleren,
het “filosoleren”.

Ze liet elk complex woord
nog niet eerder gehoord

zijn onbewust werk doen
via een visioen.

Emmy had het talent
om haar brein permanent

scherp te observeren
en te bestuderen.

Ze kon haar gedachten
meestal goed betrachten.

Ze kon haar emoties
zonder veel restricties

toestaan en verkennen
om aan ze te wennen.

Ze kon haar zintuigen
eren als getuigen.

Ze kon zich op lijden
het beste opleiden.

Ze dacht door begrijpen
haar ziel vast te grijpen.

Ze dacht zo haar ego
zijn beschermingsniveau

en zijn beperkingen
grondig te doordringen

om zichzelf te vinden
en zich te verbinden

met dit Godsbewustzijn.
Het onbegrensde zijn.

Ze zag de valstrik niet.
Hoe Eshu als bandiet

zijn rode kant liet zien,
om eens gespot nadien

ineens blauw te worden.
Het brein kon de horden,

die de trickster opriep,
waar de geest in vastliep,

nimmer overwinnen.
De duale zinnen

in het hoofd van Emmy
veroorzaakten ruzie.

Ze miste integratie,
dus verdichtte Emmy

een verhaal als kapstok,
dat pas een jas aantrok,

als een versleten ding
weer aan het haakje ging,

omdat die niet pastte
of iets haar verraste.

Dit constructivisme
leidde tot holisme.

Dat vulde elk gaatje.
Welk legpuzzelplaatje

al die losse stukjes,
als kleine constructjes,

zullen openbaren,
kon ze pas ervaren

als ze het merendeel
van het grote geheel

bij elkaar had gebracht.
Wat had ze zo bedacht?
